% generated by paper_korr.py
\begin{table}[t]
\centering
\begin{tabular}{llccc}
\hline\hline
state & $\alpha_d$ & $\alpha_d^2\langle r^2\rangle$ & dev. & pull\\
\hline
meson & 0.5 & $1.751(64)$ & $+3.8\%$ & $+1.0$\\
meson & 0.03 & $1.607(73)$ & $-4.8\%$ & $-1.1$\\
meson & 0.001 & $1.594(59)$ & $-5.6\%$ & $-1.6$\\
baryon & 0.5 & $3.855(91)$ & $+0.8\%$ & $+0.3$\\
baryon & 0.03 & $3.86(11)$ & $+0.8\%$ & $+0.3$\\
baryon & 0.001 & $3.96(13)$ & $+3.4\%$ & $+1.0$\\
\hline\hline
\end{tabular}
\caption{GFMC follows the exact Coulomb scaling law over a factor 500 in $\alpha_d$, with
$\alpha_d^2\langle r^2\rangle=27/16=1.6875$ (meson) and $3.826$ (baryon, variational).
Runs at different $\alpha_d$ are the same calculation in units of the Bohr radius, with
$d\tau\propto\alpha_d^{-2}$ and $q\propto\alpha_d$. Since a Coulomb potential has no running, $\alpha_d$ only sets
the units here, so these runs test the method at any $\alpha_d$.}
\label{tab:coulval}
\end{table}
